\scanpage{40}
\inserted{que $\operatorname{card}I<\operatorname{card}J-1$ i.e. $\operatorname{card}(J\setminus I)=\operatorname{card}(Y\setminus U)>1$.}
\deleted{On est donc dans l'un des cas suivants}

\deleted{a) $f^{-1}(I)=I$ i.e. $I$ stable par $f$.}

\deleted{b)} Donc on voit que l'on doit avoir $f^{-1}(I)=I$ i.e. $I$ stable par $f$, ce qui montre que $\tau_u$ induit en fait un \uline{automorphisme} de $Y$.
\margin{\deleted{Remplaçant \unsure{$u$}}}

\deleted{que $Y$ est une courbe elliptique}

\unsure{Idem}, en remplaçant $Y=\widehat Y\setminus I$ par $Y'=Y\cup i=\widehat Y\setminus I'$ $(I'=I-\{i\})$, on voit de $\unsure{\overline m}$ que $f$ induit un autom. de $Y'$ (cet argument marche dès que $Y'$ s'injecte dans sa jacobienne généralisée, ce qui \deleted{sauf le cas où $Y$ \unsure{est} de type $(0,\nu)$ \unsure{avec}} \inserted{est le cas grâce à l'hyp. anabélienne sur $Y$})

\deleted{\unsure{et}} donc $f$ induit l'identité sur $I$. Si le genre est $0$, on en conclut (puisque $\nu\geq3$) que $f$ est l'identité; \deleted{Reste le} Dans le cas du genre $1$, on en conclut au moins que l'image de $u$ dans $J_{\widehat Y}^0$ est égale à $1$, et on termine comme précédemment.

Ouf!
\margin{39}
